\documentclass[10pt,twocolumn]{article}

% ---------- Basic packages ----------
\usepackage[margin=1in]{geometry}
\usepackage{times}
\usepackage{amsmath, amssymb, amsfonts}
\usepackage{graphicx}
\usepackage{subcaption}
\usepackage{booktabs}
\usepackage{multirow}
\usepackage{siunitx}
\usepackage{url}
\usepackage[colorlinks=true,allcolors=blue]{hyperref}
\usepackage[round,authoryear]{natbib}
\usepackage{microtype}
\usepackage{setspace}
\usepackage{caption}
\usepackage{tabularx}

% Figures folder (compile from repo root: pdflatex report/main.tex)
\graphicspath{{C:/UT Austin MS CS/CML/PromiseAwareETA/figures/}}

% 1.5 line spacing to match course requirements
\onehalfspacing

% Avoid stretched pages that can cause ugly spacing
\raggedbottom

% ---------- Metadata ----------
\title{Promise-Aware ETA: Calibrated Delivery-Time Uncertainty and Offline Delivery-Promise Optimization on the Olist E-Commerce Dataset}

\author{
  Your Name\thanks{Affiliation, address, email.}
}

\date{\today}

\begin{document}
\maketitle

\begin{abstract}
Accurate and trustworthy delivery-time promises are now a key part of the online shopping experience. In practice, many e-commerce systems still make a single point estimate of delivery time and then add a hand-tuned buffer. This often leads to a mix of broken promises (late deliveries) and unnecessarily long quoted lead times. 

This report introduces \emph{Promise-Aware ETA}, an end-to-end pipeline that (i) models order-level delivery-time uncertainty using quantile regression, (ii) calibrates these quantiles with split conformal prediction, and (iii) evaluates different delivery-promise policies offline using an asymmetric cost function and fairness diagnostics across sellers and regions. The approach is instantiated on the Brazilian E-Commerce Public Dataset by Olist, a multi-table dataset with roughly 100\,000 orders.

Using LightGBM-based quantile regression models combined with conformalized quantile regression, the final model achieves 90\,\% nominal prediction intervals with empirical coverage of 0.8997 and a mean interval width of 22.0 days on a 2018 hold-out set. A simple promise policy that chooses the calibrated 0.9 quantile as the promised date increases coverage from 51.1\,\% to 84.2\,\% compared to a median-based policy, while reducing the expected asymmetric cost from 22.35 to 15.44 penalty-weighted days per order under a 5:1 late-versus-early penalty ratio. Regional fairness analysis shows coverage ranging from 0.60 in the sparsely sampled North region to 0.87 in the Central-West region, highlighting data sparsity and heterogeneity as ongoing risks. At the seller level, coverage averages 0.86 but exhibits a heavy-tailed distribution driven by long-tail sellers with few orders.

Beyond the empirical results, the main contribution is a modular, leakage-safe pipeline for feature engineering, calibration, and policy evaluation, supported by a reproducibility checklist and experiment logs. The report closes with implications for e-commerce logistics and possible extensions to richer policy classes, fairness-aware objectives, and production deployment.
\end{abstract}

\section{Introduction}
Delivery-time expectations have become a central competitive factor for e-commerce platforms. On most marketplaces, the \emph{promised delivery date} is visible on the product page and at checkout. It affects whether a customer completes a purchase, how they feel about the order, and whether they return. When deliveries are late, platforms lose trust, incur customer-support costs, and may need to issue refunds or compensation. When promises are overly conservative, customers perceive long lead times and may choose a competitor instead.

In production, many systems still use a simple pattern: estimate a single expected delivery lag (using a regression model or lookup table) and then add a heuristic buffer such as ``+2 days'' or ``max(expected, carrier SLA).'' This approach does not handle heavy-tailed, highly variable delays caused by geography, seller behavior, or random disruptions in the logistics network. It also gives product and operations teams little control over the trade-off between the risk of being late and the length of the promise.

Recent work in forecasting and uncertainty quantification argues for modeling the \emph{full conditional distribution} of delivery times, rather than only a mean or median \citep{koenker1978,koenker2001,lightgbm2017,zhang2023calibration}. Quantile regression and conformal prediction fit well in this setting: quantile models estimate conditional quantiles $Q_\tau(Y \mid X)$ of the delivery lag $Y$ given features $X$, and conformal prediction then calibrates the resulting intervals to achieve finite-sample coverage guarantees on held-out data \citep{romano2019conformalized,angelopoulos2021gentle}.

This report takes a decision-making viewpoint. Instead of asking only ``How accurate is the ETA?'', the core question is: given calibrated predictive uncertainty, which promise policy minimizes a business-relevant asymmetric cost while satisfying reasonable fairness and reliability constraints? We study this question on the Olist dataset \citep{olistdataset} using a pipeline we refer to as \emph{Promise-Aware ETA}.

\paragraph{Research questions.}
Concretely, the report focuses on three questions:
\begin{enumerate}
  \item \textbf{Calibration.} After training quantile models for order-level delivery lag and calibrating them with split conformal prediction, how close can we get to the desired coverage, and with what interval widths?
  \item \textbf{Policy optimality.} Within a family of quantile-based promise policies (for example, promising at the calibrated 0.5 or 0.9 quantile), which choices minimize an asymmetric cost function that penalizes late deliveries more heavily than early ones?
  \item \textbf{Fairness and robustness.} How do these policies behave across key segments such as sellers and regions? Are there systematic coverage or cost gaps that would be problematic for a real-world deployment?
\end{enumerate}

\paragraph{Contributions.}
The main contributions of this work are:
\begin{itemize}
  \item A leakage-safe feature engineering and modeling pipeline for the Olist dataset, including rolling seller dispatch statistics and geographic features.
  \item A conformalized quantile regression (CQR) setup that wraps base quantile models to produce calibrated prediction intervals at a chosen target coverage.
  \item An offline promise-policy simulator with an asymmetric cost function and group-wise fairness diagnostics across regions and sellers.
  \item Empirical evidence that a calibrated 0.9-quantile promise policy can both reduce expected cost and substantially improve coverage relative to a median-based policy.
  \item A reproducibility checklist and experiment logs (configuration, metrics, and artifacts) that capture the modeling process end-to-end.
\end{itemize}

\section{Related Work}
\subsection{Delivery-time prediction and e-commerce logistics}
The Brazilian E-Commerce Public Dataset by Olist is widely used as a benchmark for marketplace behavior, logistics, and customer satisfaction studies \citep{olistdataset}. Prior work has applied statistical and machine learning models to predict delivery times from order, seller, and product-level features \citep{brunilogistics2023,rokoss2024deliverytime}. Case studies from large retailers show that improved ETA systems can significantly impact last-mile logistics performance and customer satisfaction \citep{walmart2019cdt,kardinal2025lastmile}.

Alongside prediction, there is a growing body of work on optimization and planning in last-mile logistics. Machine-learning-enhanced heuristics have been proposed for vehicle routing and delivery-time constraints, often outperforming classical operations research baselines \citep{brunilogistics2023}. However, many of these systems still rely on point ETAs with fixed or ad-hoc buffers rather than explicit uncertainty estimates.

\subsection{Quantile regression and pinball loss}
Quantile regression extends ordinary least squares by directly modeling conditional quantiles instead of the conditional mean \citep{koenker1978,koenker2001}. The standard approach minimizes the \emph{pinball loss} (or check loss), an asymmetric linear loss whose asymmetry parameter selects which quantile is being estimated. There are many extensions, including quantile random forests, gradient-boosted quantile models, and formulations that jointly train multiple quantiles via expected pinball loss \citep{biau2024pinball,tagasovska2019epinball}.

For tabular data, gradient boosting decision trees---particularly LightGBM---have become standard tools for quantile regression because they tend to perform well and train efficiently \citep{lightgbm2017}. Their ability to handle mixed feature types is especially useful for structured e-commerce data such as orders, sellers, and locations.

\subsection{Conformal prediction and calibration}
Conformal prediction provides a general, distribution-free framework for constructing prediction sets with finite-sample coverage guarantees under exchangeability \citep{vovk2005algorithmic,angelopoulos2021gentle}. Conformalized quantile regression (CQR) adapts these ideas to quantile models by adjusting lower and upper quantiles using calibration residuals, so that the resulting prediction intervals achieve the desired coverage \citep{romano2019conformalized}. Recent work has introduced adaptive and shift-aware variants of CQR for tabular data \citep{sousa2024adaptive,zhang2023calibration}.

In parallel, there is extensive work on calibration metrics and reliability diagrams for probabilistic models, including coverage-based measures and calibration error metrics \citep{guo2017calibration,levi2022uncertainty}. In this report, we focus primarily on coverage, violation rates, and interval width as our main calibration diagnostics.

\subsection{Fairness and offline policy evaluation}
Research on algorithmic fairness has introduced several group fairness notions and has shown that calibration and parity across groups often conflict \citep{barocas2019fairness,pleiss2017fairness,kleinberg2017impossibility}. For logistics and on-demand platforms, fairness questions arise in routing and work allocation for delivery workers and riders \citep{gupta2022fairfoody,martinez2024fairrouting}.

Evaluating new decision policies from historical logs, without running online experiments, is a key challenge in contextual bandits and reinforcement learning. A variety of importance-weighted and doubly robust estimators have been proposed for off-policy evaluation \citep{dudik2011doubly,jiang2016doubly}. The offline promise simulator in this report is simpler: it assumes that the new promise policy does not influence the actual delivery lag, so we can treat realized lags as fixed outcomes.

\section{Materials and Data Sources}
\subsection{Dataset}
The Brazilian E-Commerce Public Dataset by Olist is a public Kaggle dataset with around 100\,000 orders split across multiple relational tables \citep{olistdataset}. The key tables used in this work are:
\begin{itemize}
  \item \textbf{Orders:} order status and timestamps (purchase, approval, shipment, and delivery dates).
  \item \textbf{Order items:} product category, seller ID, freight value, and shipping-limit dates.
  \item \textbf{Customers and sellers:} location data, aggregated into Brazilian macro-regions.
\end{itemize}

\subsection{Target definition}
The target variable is the approval-to-delivery lag in days:
\begin{equation}
  Y = \max\bigl\{0,\; t_{\text{delivered}} - t_{\text{approved}}\bigr\},
\end{equation}
where $t_{\text{delivered}}$ and $t_{\text{approved}}$ are the delivery and approval timestamps, respectively. Orders with missing or inconsistent timestamps (for example, delivery recorded before approval) are filtered out.

\subsection{Feature engineering}
Feature generation follows a structured pipeline. The main feature families are:
\begin{itemize}
  \item \textbf{Order-level features:} freight value, number of items, product category encodings, payment method, installment count, and a shipping-limit slack feature (days between approval and the shipping deadline).
  \item \textbf{Geographical features:} customer--seller distance computed from latitude and longitude via the Haversine formula, seller region, and an indicator for intra-region shipments.
  \item \textbf{Temporal features:} purchase month, day of week, and proximity to major holidays.
  \item \textbf{Seller dispatch history:} rolling 30/60/90-day per-seller statistics, including median dispatch lag, 90th percentile dispatch lag, share of orders shipped before the shipping deadline, and the difference between recent and baseline medians. These statistics are computed in a leakage-safe way by only using orders with approval timestamps strictly earlier than the target order.
\end{itemize}

\subsection{Train--validation split}
To better approximate a real deployment scenario, we use a temporal split:
\begin{itemize}
  \item Training: orders approved between 2017-01-01 and 2017-12-31.
  \item Validation: orders approved between 2018-01-01 and 2018-06-30.
\end{itemize}

\section{Methods}
\subsection{Problem formulation}
For each order $i$, let $X_i \in \mathbb{R}^d$ be the feature vector and $Y_i \in \mathbb{R}_{\ge 0}$ the realized delivery lag. Rather than modeling $\mathbb{E}[Y \mid X]$, our goal is to estimate conditional quantiles $Q_\tau(X)$ for $\tau \in \{0.1, 0.5, 0.9\}$. The base quantile models minimize the pinball loss:
\begin{equation}
  \ell_\tau(y,\hat{q}) = 
  \begin{cases}
    \tau (y - \hat{q}), & y \ge \hat{q},\\[2pt]
    (\tau - 1)(y - \hat{q}), & y < \hat{q}.
  \end{cases}
\end{equation}

\subsection{Base quantile models}
We consider three quantile model families:
\begin{enumerate}
  \item \textbf{Linear quantile regression}, trained with a smoothed pinball loss and $\ell_2$ regularization.
  \item \textbf{Histogram-based gradient boosting (HGB)}, using a gradient boosting regressor configured for quantile loss.
  \item \textbf{LightGBM quantile regression}, using histogram-based gradient boosted trees \citep{lightgbm2017}.
\end{enumerate}

Figure~\ref{fig:f1-pinball} shows the pinball loss of the baseline models across quantiles on a synthetic benchmark, which we use to validate the training and evaluation pipeline before applying it to the Olist data.

\begin{figure}[t]
  \centering
  \includegraphics[width=0.75\linewidth]{F1_pinball_by_tau.png}
  \caption{Pinball loss by quantile $\tau$ for baseline models on a synthetic benchmark dataset. Lower values indicate better performance.}
  \label{fig:f1-pinball}
\end{figure}

\subsection{Split-conformal calibration}
Base quantile models are usually not perfectly calibrated: the empirical coverage of nominal $(1-\alpha)$ intervals can deviate noticeably from the target. We therefore apply split conformal calibration following CQR \citep{romano2019conformalized}.

Let $\hat{Q}_{\alpha}^{\text{base}}(x)$ and $\hat{Q}_{1-\alpha}^{\text{base}}(x)$ denote the lower and upper base quantiles for target coverage $1-\alpha$. We split the training data into a proper training set and a separate calibration set. On the calibration set, we compute non-conformity scores
\begin{equation}
  s_i = \max\left\{\hat{Q}_{\alpha}^{\text{base}}(X_i) - Y_i,\; Y_i - \hat{Q}_{1-\alpha}^{\text{base}}(X_i),\; 0\right\}.
\end{equation}
We then take $\hat{\Delta}$ to be the empirical $(1-\alpha)$ quantile of $\{s_i\}$ and form calibrated intervals
\begin{equation}
  [\hat{L}(x),\hat{U}(x)] = 
  \big[\hat{Q}_{\alpha}^{\text{base}}(x) - \hat{\Delta},\; \hat{Q}_{1-\alpha}^{\text{base}}(x) + \hat{\Delta}\big].
\end{equation}

Figure~\ref{fig:f2-reliability} compares nominal and empirical coverage across target levels, and Figure~\ref{fig:f3-width-coverage} reports the corresponding average interval widths.

\begin{figure}[t]
  \centering
  \includegraphics[width=0.75\linewidth]{F2_reliability.png}
  \caption{Reliability diagram: empirical coverage vs.\ nominal target coverage for the calibrated model on the Olist validation set. The dashed line shows perfect calibration.}
  \label{fig:f2-reliability}
\end{figure}

\begin{figure}[t]
  \centering
  \includegraphics[width=0.75\linewidth]{F3_width_vs_coverage.png}
  \caption{Mean prediction interval width (in days) as a function of target coverage for the conformalized quantile model.}
  \label{fig:f3-width-coverage}
\end{figure}

\subsection{Promise policies and asymmetric cost}
Once we have calibrated quantiles and prediction intervals, we define a \emph{promise policy} as a mapping from features $X$ to a promised lag $P(X)$ in days. We consider simple quantile-based policies of the form
\begin{equation}
  P_\tau(X) = \hat{Q}_{\tau}^{\text{cal}}(X),
\end{equation}
where $\hat{Q}_{\tau}^{\text{cal}}$ is the calibrated quantile function. We focus on two policies:
\begin{itemize}
  \item \textbf{p50:} promise at the calibrated median ($\tau = 0.5$).
  \item \textbf{p90:} promise at the calibrated upper quantile ($\tau = 0.9$).
\end{itemize}

To compare policies, we use an asymmetric cost function that penalizes late deliveries more than early deliveries. Let $p_i$ be the promise lag and $y_i$ the realized lag for order $i$. The cost is
\begin{equation}
  C(y_i, p_i) =
  \begin{cases}
    \lambda_{\text{late}} (y_i - p_i), & y_i > p_i,\\[2pt]
    \lambda_{\text{early}} (p_i - y_i), & y_i < p_i,\\[2pt]
    0, & y_i = p_i,
  \end{cases}
\end{equation}
with default penalties $\lambda_{\text{late}} = 5$ and $\lambda_{\text{early}} = 1$.

Figure~\ref{fig:f4-cost-tau} shows the expected cost as a function of the promised quantile $\tau$. Figure~\ref{fig:f5-lpr-apd} illustrates the trade-off between late-promise rate (LPR) and average promise days (APD) for the same family of policies.

\begin{figure}[t]
  \centering
  \includegraphics[width=0.75\linewidth]{F4_cost_vs_tau.png}
  \caption{Expected asymmetric cost as a function of promise quantile $\tau$ for quantile-based policies.}
  \label{fig:f4-cost-tau}
\end{figure}

\begin{figure}[t]
  \centering
  \includegraphics[width=0.75\linewidth]{F5_lpr_apd_frontier.png}
  \caption{Late-promise rate (LPR) vs.\ average promise days (APD) frontier for different quantile-based policies. The recommended operating point at $\tau=0.9$ balances coverage and promise length.}
  \label{fig:f5-lpr-apd}
\end{figure}

\subsection{Fairness metrics}
We assess fairness across two types of groups:
\begin{itemize}
  \item \textbf{Regions:} seller macro-regions (North, Northeast, Central-West, Southeast, South).
  \item \textbf{Sellers:} individual seller IDs.
\end{itemize}
For each group $g$ and policy $P$, we compute group coverage, expected cost, and their differences or ratios relative to reference groups.

\section{Results}
\subsection{Calibration performance}
On the 2018 validation window, the LightGBM + CQR model achieves:
\begin{itemize}
  \item Target coverage: 90\,\%.
  \item Observed coverage: 0.8997.
  \item Mean interval width: 22.03 days.
\end{itemize}
These metrics are visualized in Figures~\ref{fig:f2-reliability} and \ref{fig:f3-width-coverage}.

\subsection{Promise-policy comparison}
Table~\ref{tab:policies} compares the p50 and p90 policies on the validation set.

\begin{table}[t]
  \centering
  \caption{Policy metrics for two quantile-based promise policies on the Olist validation set. The cost uses $\lambda_{\text{late}}=5$ and $\lambda_{\text{early}}=1$.}
  \label{tab:policies}
  \setlength{\tabcolsep}{4pt}
  \renewcommand{\arraystretch}{1.1}
  \resizebox{\linewidth}{!}{%
    \begin{tabular}{lccccccc}
      \toprule
      Policy & Quantile & Coverage & Avg.\ P. & Avg.\ A. & Avg.\ Late & Avg.\ Early & Exp.\ Cost \\[1pt]
      \midrule
      p50 & 0.5 & 0.511 & 11.02 & 13.33 & 4.11 & 1.80 & 22.35 \\
      p90 & 0.9 & 0.842 & 19.98 & 13.33 & 1.47 & 8.12 & 15.44 \\
      \bottomrule
    \end{tabular}%
  }
\end{table}

Figure~\ref{fig:f5-lpr-apd} shows that moving from a median-based policy to a higher quantile sharply reduces the late-promise rate, at the cost of longer promised lead times. Given that late days are penalized more heavily, the p90 policy ends up with a lower expected cost than p50, despite the longer promises.

\subsection{Segment-wise policies}
To capture heterogeneity, we also consider segment-specific policies, where the chosen quantile $\tau$ varies by distance and historical dispatch-delay segments. Figure~\ref{fig:f6a-segment-heatmap} shows the expected cost across segments, annotated with the selected $\tau$, and Figure~\ref{fig:f6b-tau-frequency} summarizes how often each quantile is used.

\begin{figure}[t]
  \centering
  \includegraphics[width=0.75\linewidth]{F6a_expected_cost_heatmap_annotated_tau.png}
  \caption{Segmented expected cost heatmap (distance decile vs.\ dispatch-delay decile) for a segment-wise policy, annotated with the promise quantile $\tau$ in each segment.}
  \label{fig:f6a-segment-heatmap}
\end{figure}

\begin{figure}[t]
  \centering
  \includegraphics[width=0.6\linewidth]{F6b_tau_frequency.png}
  \caption{Frequency of selected promise quantiles $\tau$ across segments in the segment-wise policy.}
  \label{fig:f6b-tau-frequency}
\end{figure}

\subsection{Regional fairness}
Figure~\ref{fig:f7-region-coverage} and Table~\ref{tab:regions} summarize regional coverage and cost for the p90 policy.

\begin{table}[t]
  \centering
  \caption{Regional fairness metrics for the p90 policy. Coverage and expected cost are reported along with differences and ratios to reference regions.}
  \label{tab:regions}
  \setlength{\tabcolsep}{4pt}
  \renewcommand{\arraystretch}{1.1}
  \resizebox{\linewidth}{!}{%
    \begin{tabular}{lrrrr}
      \toprule
      Region & Coverage & $\Delta$Cov vs.\ CW & Exp.\ Cost & Cost Ratio vs.\ NE \\[1pt]
      \midrule
      Central-West & 0.873 & 0.000 & 14.69 & 0.88 \\
      South        & 0.850 & -0.023 & 15.37 & 0.92 \\
      Northeast    & 0.848 & -0.025 & 16.73 & 1.00 \\
      Southeast    & 0.841 & -0.032 & 15.44 & 0.96 \\
      North        & 0.600 & -0.273 & 11.62 & 0.69 \\
      \bottomrule
    \end{tabular}%
  }
\end{table}

\begin{figure}[t]
  \centering
  \includegraphics[width=0.75\linewidth]{F7_region_coverage.png}
  \caption{Regional coverage for the $\tau=0.9$ promise policy. Bars show group coverage; the dashed line marks the Central-West reference region.}
  \label{fig:f7-region-coverage}
\end{figure}

Four regions cluster around 0.84--0.87 coverage, with cost ratios within roughly $\pm 12\,\%$ of the Northeast reference. The North region stands out with 0.60 coverage, driven by a relatively small number of orders and potentially different logistics conditions.

\subsection{Seller-level fairness}
At the seller level (details not fully tabulated here), we observe:
\begin{itemize}
  \item Mean coverage of 0.86 and median coverage of 0.93.
  \item An interquartile range of 0.80--1.00.
  \item A long tail of low-coverage sellers, largely overlapping with sellers that have few orders.
\end{itemize}
Figure~\ref{fig:f8-seller-coverage} shows the distribution of seller-level coverage.

\begin{figure}[t]
  \centering
  \includegraphics[width=0.75\linewidth]{F8_seller_coverage_hist.png}
  \caption{Distribution of seller-level coverage for the $\tau=0.9$ policy. The vertical line marks the global 0.9 target.}
  \label{fig:f8-seller-coverage}
\end{figure}

\subsection{Feature contributions and diagnostics}
Figure~\ref{fig:f9-feature-importance} shows gain-based feature importance for the LightGBM quantile model at $\tau=0.9$, emphasizing the role of seller dispatch history and geographical features.

\begin{figure}[t]
  \centering
  \includegraphics[width=0.75\linewidth]{F9_feature_importance_tau90.png}
  \caption{Feature importance (gain-based) for the LightGBM quantile model at $\tau=0.90$.}
  \label{fig:f9-feature-importance}
\end{figure}

Figure~\ref{fig:f10-distance-slack} plots promise slack (promise minus actual) against corridor distance, making it easier to see where late orders occur.

\begin{figure}[t]
  \centering
  \includegraphics[width=0.75\linewidth]{F10_violation_vs_distance.png}
  \caption{Slack (promise minus actual delivery time) vs.\ corridor distance for the $\tau=0.9$ policy; late orders are highlighted.}
  \label{fig:f10-distance-slack}
\end{figure}

Finally, Figure~\ref{fig:f11-temporal-coverage} shows coverage over time, and Figure~\ref{fig:f12-ablation} summarizes ablation experiments.

\begin{figure}[t]
  \centering
  \includegraphics[width=0.75\linewidth]{F11_coverage_over_time.png}
  \caption{Empirical coverage of the $\tau=0.9$ policy over time (weekly buckets). The horizontal line shows the nominal 0.9 target.}
  \label{fig:f11-temporal-coverage}
\end{figure}

\begin{figure}[t]
  \centering
  \includegraphics[width=0.75\linewidth]{F12_ablation_tornado.png}
  \caption{Ablation tornado plot: change in mean pinball loss when removing feature families, compared to the full model.}
  \label{fig:f12-ablation}
\end{figure}

\section{Discussion}
\subsection{Interpretation}
The main takeaway is that promise policies based on calibrated quantile predictions can substantially improve both reliability and cost compared to naive median-based promises. A 0.9-quantile policy delivers much higher coverage and lower expected cost under a realistic asymmetric penalty that heavily discourages late deliveries.

The calibrated intervals are relatively wide, which reflects real variability in the Olist data, including long-distance and cross-region shipments. In a live system, additional real-time features (for example, day-of-route information or carrier-level signals) would likely narrow these intervals and sharpen the trade-offs.

\subsection{Fairness implications}
The regional and seller-level diagnostics show that having good average calibration is not enough. Some segments, notably the sparsely sampled North region, see much lower coverage. Potential mitigation strategies include segment-wise calibration, minimum coverage constraints by group, and hierarchical models that borrow statistical strength across regions \citep{pleiss2017fairness,pfohl2021fairrisk}.

\subsection{Limitations}
Key limitations of this work include:
\begin{itemize}
  \item The Olist dataset is static and lacks real-time operational signals commonly used in modern ETA systems.
  \item The policy class is intentionally simple, restricted to quantile-based rules.
  \item Conformal prediction guarantees marginal coverage, not group-wise coverage across all subpopulations.
  \item The offline evaluation assumes that the promise does not affect the realized delivery lag, which may not be fully realistic if promises influence operational behavior.
\end{itemize}

\section{Conclusion and Future Work}
This report has presented Promise-Aware ETA, a reproducible pipeline for modeling calibrated delivery-time uncertainty and optimizing delivery promises offline on the Olist dataset. LightGBM-based quantile models, calibrated with split conformal prediction, achieve near-nominal coverage. Combined with quantile-based promise policies and an asymmetric cost function, these models yield promise strategies that are more reliable and more cost-effective than naive median-based baselines. Fairness diagnostics reveal heterogeneous performance across regions and sellers, underscoring the need for segment-aware calibration and policy design.

Future directions include broadening the policy class to richer, feature-dependent rules; incorporating fairness constraints directly into policy optimization; and integrating real-time signals in an online evaluation or A/B testing framework.

\section*{Acknowledgments}
[Optional: acknowledge advisors, collaborators, or funding sources.]

\bibliographystyle{plainnat}
\begin{thebibliography}{99}

\bibitem[Angelopoulos and Bates(2021)]{angelopoulos2021gentle}
A.~N. Angelopoulos and S.~Bates.
\newblock A gentle introduction to conformal prediction and distribution-free uncertainty quantification.
\newblock \emph{arXiv preprint arXiv:2107.07511}, 2021.

\bibitem[Barocas et~al.(2019)]{barocas2019fairness}
S.~Barocas, M.~Hardt, and A.~Narayanan.
\newblock \emph{Fairness and Machine Learning}.
\newblock MIT Press, 2019.

\bibitem[Biau and B{\"o}hning(2024)]{biau2024pinball}
G.~Biau and D.~B{\"o}hning.
\newblock Pinball boosting of regression quantiles.
\newblock \emph{Computational Statistics \& Data Analysis}, 2024.

\bibitem[Bruni et~al.(2023)]{brunilogistics2023}
M.~E. Bruni et~al.
\newblock A machine learning optimization approach for last-mile delivery with time windows.
\newblock \emph{Computers \& Operations Research}, 2023.

\bibitem[Dud{\'\i}k et~al.(2011)]{dudik2011doubly}
M.~Dud{\'\i}k, J.~Langford, and L.~Li.
\newblock Doubly robust policy evaluation and learning.
\newblock In \emph{Proceedings of ICML}, 2011.

\bibitem[Guo et~al.(2017)]{guo2017calibration}
C.~Guo, G.~Pleiss, Y.~Sun, and K.~Q. Weinberger.
\newblock On calibration of modern neural networks.
\newblock In \emph{Proceedings of ICML}, 2017.

\bibitem[Gupta et~al.(2022)]{gupta2022fairfoody}
A.~Gupta et~al.
\newblock FairFoody: Bringing in fairness in food delivery.
\newblock In \emph{Proceedings of The Web Conference}, 2022.

\bibitem[Jiang and Li(2016)]{jiang2016doubly}
N.~Jiang and L.~Li}.
\newblock Doubly robust off-policy value evaluation for reinforcement learning.
\newblock In \emph{Proceedings of ICML}, 2016.

\bibitem[Kardinal(2025)]{kardinal2025lastmile}
Kardinal.
\newblock Machine learning: 8 innovations for last-mile delivery.
\newblock Technical report, 2025.

\bibitem[Kleinberg et~al.(2017)]{kleinberg2017impossibility}
J.~Kleinberg, S.~Mullainathan, and M.~Raghavan.
\newblock Inherent trade-offs in the fair determination of risk scores.
\newblock In \emph{Proceedings of ITCS}, 2017.

\bibitem[Koenker(2001)]{koenker2001}
R.~Koenker.
\newblock Quantile regression.
\newblock \emph{Journal of Economic Perspectives}, 15(4):143--156, 2001.

\bibitem[Koenker and Bassett(1978)]{koenker1978}
R.~Koenker and G.~Bassett.
\newblock Regression quantiles.
\newblock \emph{Econometrica}, 46(1):33--50, 1978.

\bibitem[Levi et~al.(2022)]{levi2022uncertainty}
D.~Levi et~al.
\newblock Evaluating and calibrating uncertainty prediction in deep neural networks.
\newblock \emph{International Journal of Computer Vision}, 130, 2022.

\bibitem[Ke et~al.(2017)]{lightgbm2017}
G.~Ke et~al.
\newblock LightGBM: A highly efficient gradient boosting decision tree.
\newblock In \emph{Advances in Neural Information Processing Systems}, 2017.

\bibitem[Mart{\'\i}nez-Sykora et~al.(2024)]{martinez2024fairrouting}
A.~Mart{\'\i}nez-Sykora et~al.
\newblock Exploring fairness in routing and scheduling for food delivery platforms.
\newblock \emph{Expert Systems with Applications}, 2024.

\bibitem[Olist(2018)]{olistdataset}
Olist.
\newblock Brazilian E-Commerce Public Dataset by Olist.
\newblock Kaggle data set, 2018.

\bibitem[Pfohl et~al.(2021)]{pfohl2021fairrisk}
S.~R. Pfohl et~al.
\newblock An empirical characterization of fair machine learning for clinical risk prediction.
\newblock \emph{Journal of Biomedical Informatics}, 113, 2021.

\bibitem[Pleiss et~al.(2017)]{pleiss2017fairness}
G.~Pleiss et~al.
\newblock On fairness and calibration.
\newblock In \emph{Advances in Neural Information Processing Systems}, 2017.

\bibitem[Rokoss et~al.(2024)]{rokoss2024deliverytime}
A.~Rokoss et~al.
\newblock Case study on delivery time determination using a machine learning approach.
\newblock \emph{Journal of Intelligent Manufacturing}, 2024.

\bibitem[Romano et~al.(2019)]{romano2019conformalized}
Y.~Romano, E.~Patterson, and E.~J. Cand{\`e}s.
\newblock Conformalized quantile regression.
\newblock In \emph{Advances in Neural Information Processing Systems}, 2019.

\bibitem[Sousa et~al.(2024)]{sousa2024adaptive}
M.~Sousa et~al.
\newblock Improving conformalized quantile regression through adaptive calibration.
\newblock \emph{Expert Systems with Applications}, 2024.

\bibitem[Tagasovska and Lopez-Paz(2019)]{tagasovska2019epinball}
N.~Tagasovska and D.~Lopez-Paz.
\newblock Single-model uncertainties for deep learning.
\newblock In \emph{Advances in Neural Information Processing Systems}, 2019.

\bibitem[Vovk et~al.(2005)]{vovk2005algorithmic}
V.~Vovk, A.~Gammerman, and G.~Shafer.
\newblock \emph{Algorithmic Learning in a Random World}.
\newblock Springer, 2005.

\bibitem[Walmart Global Tech(2019)]{walmart2019cdt}
Walmart Global Tech.
\newblock Customer delivery time prediction using machine learning.
\newblock Technical report, 2019.

\bibitem[Zhang et~al.(2023)]{zhang2023calibration}
X.~Zhang et~al.
\newblock Calibration in machine learning uncertainty quantification.
\newblock \emph{arXiv preprint arXiv:2309.06240}, 2023.

\end{thebibliography}

\end{document}
